\clearpage
\section*{Question 3}
\addcontentsline{toc}{section}{Question 3}

\textbf{Problem Statement:} Calculate the initial conditions that would correspond to an idle steady-state operation with $T_m = 0$ and excitation $E_{xfd} = 1.0$ pu, and set the initial conditions appropriately in your model. Implement the following transient study:

\begin{itemize}
    \item At $t = 1$ s, the excitation $E_{xfd}$ is stepped down to 0.8 pu.
    \item Then, at $t = 7$ s, the excitation $E_{xfd}$ is stepped up to 1.2 pu.
    \item Then, at $t = 15$ s, the mechanical torque is stepped to $T_m = 0.85$ pu supplying power from the mechanical system to the electrical system, and the model is continued to run till $t = 25$ s.
\end{itemize}

Plot variables $i_{as}$, $i_{fd}$, $i_{kq}$, $T_e$, $\omega_r$, $\theta_r$ (in electrical degrees), as well as real power $P_e$ and reactive power $Q_e$. Based on the transients observed, briefly explain the effect of the excitation and torque on the active and reactive power and the type of transients seen.

\vspace{1em}
\noindent\rule{\textwidth}{0.4pt}
\vspace{1em}

\textbf{Solution:}

\subsection*{Initial Conditions Calculation}

% TODO: Show calculation of steady-state initial conditions

\subsection*{Simulation Results}

% TODO: Add plots of all required variables

\subsection*{Analysis of Excitation Effects}

% TODO: Explain effect of excitation changes on Pe and Qe

\subsection*{Analysis of Torque Effects}

% TODO: Explain effect of mechanical torque on system behavior

\vspace{1em}
\noindent\rule{\textwidth}{0.4pt}
